\Askhsh{20712}{4}
\begin{ylh}{1}
\item 3.4 Οι τριγωνομετρικές συναρτήσεις
\item 3.5 Βασικές τριγωνομετρικές εξισώσεις
\end{ylh}
Σε μια θαλάσσια περιοχή, λόγω της παλίρροιας, η στάθμη των υδάτων αυξομειώνεται. Στο παρακάτω σχήμα δίνεται η γραφική παράσταση της ημιτονοειδούς συνάρτησης $f$, που δίνει σε μέτρα το ύψος της στάθμης των υδάτων συναρτήσει του χρόνου $t$ σε ώρες.
\begin{center}
\begin{tikzpicture}[scale=0.5, >=latex]
    \draw[very thin, color=gray!30] (-1,-1) grid (26,5);
    \draw[->] (-1,0) -- (27,0) node[right] {$t$ (h)};
    \draw[->] (0,-1) -- (0,6) node[above] {$y$ (m)};
    \foreach \x in {3,6,9,12,15,18,21,24}
        \draw (\x,5pt) -- (\x,-5pt) node[below, font=\small] {\x};
    \foreach \y in {1,2,3,4,5}
        \draw (5pt,\y) -- (-5pt,\y) node[left, font=\small] {\y};
    \node[below left, font=\small] at (0,0) {$0$};
    % f(t) = 1.5*sin(pi/6 * t) + 1.5
    % Peaks at 3, 15... (y=3)
    % Troughs at 9, 21... (y=0)
    \draw[very thick, color=maincolor, samples=200, smooth, domain=0:26] plot (\x, {1.5*sin(deg(pi/6 * \x)) + 1.5});
    % Levels
    \draw[dashed] (0,3) -- (26,3) node[right] {$3m$};
    \draw[dashed] (0,1.5) -- (26,1.5) node[right, font=\small] {$1,5m$};
\end{tikzpicture}
\end{center}
\begin{alist}
\item την υψομετρική διαφορά ανάμεσα στην υψηλότερη στάθμη (πλημμυρίδα) και τη χαμηλότερη στάθμη (άμπωτη).
\monades{6}
\item την περίοδο του φαινομένου της παλίρροιας.
\monades{6}
\item τον τύπο της συνάρτησης $f$.
\monades{6}
\item ποιες ώρες, στη διάρκεια μιας ημέρας, η στάθμη των υδάτων είναι $\frac{3}{2}$ μέτρα.
\monades{7}
\end{alist}

